\section{Cook--Reckhow: encoding, comparisons and examples}

\begin{annexitem}{annex:p7.cr.map-relation}{thm:p07-cr-map-relation}{Formula-in-code conversion without output search}
\begin{proof}
Given f, compute it on x and compare its resulting string with y.
This takes polynomial time in \(|x|+|y|\), including reading and
comparing the output. Its existential range is L since f is onto.
Conversely, parse the paired code. If parsing fails return \(y_0\).
Otherwise run the relation decider on the explicitly stored w and y.
Its time bound is polynomial in the paired code length, and output
y has length at most that length. If the decider rejects return
the fixed \(y_0\). All outputs lie in L; every y in L has a
witness and hence a paired proof mapping to it. The map is total,
polynomial-time and onto. It does not try to reconstruct y by
enumerating candidate formulas associated with w.

If each y has a witness w of length at most q(|y|), then its paired
proof has length \(O(|y|+q(|y|)+1)\), a polynomial bound.
Conversely, a short paired proof of a y different from \(y_0\)
contains a relation witness of no greater length than its code.
For \(y_0\) fix one witness once and absorb its constant length
in the bounding polynomial. For a map's graph relation, its proof
string itself is the witness. These arguments prove boundedness
equivalence in both directions, including the default case.
\end{proof}
\end{annexitem}

\begin{annexitem}{annex:p7.cr.np-languages}{thm:p07-cr-np-languages}{Both directions of source Proposition 1.4}
\begin{proof}
If L is empty, a machine rejecting every input recognizes it.
For nonempty \(L\in\mathsf{NP}\), choose a nondeterministic
machine M recognizing it within \(t(|y|)\) steps for a polynomial t.
A certificate is a complete accepting computation, together with
its input y. One can check each successive machine configuration
and the final accepting state in polynomial time in that code:
its state, head location, input and work tape contents are explicit,
and every transition is from a fixed finite transition table.
There are \(t(|y|)+1\) configurations, each of polynomial size,
so an accepting input has a polynomial-length certificate.
Map a valid accepting-computation code to its stored y and any
invalid string to a fixed \(y_0\in L\). The preceding encoding
lemma proves this a total polynomial-time onto proof map.

Conversely, let f be polynomial-time and onto, with proof bound q.
On input y nondeterministically guess a string of at most q(|y|)
symbols, compute f on that string, and accept just when the result
equals y. The verification time is polynomial in |y| by composition
of the polynomial proof-length and computation bounds. Onto
boundedness supplies an accepting branch for each y in L, and
the range restriction excludes acceptance outside L. This is an
NP recognizer. No deterministic procedure for finding the short
certificate was obtained in either direction.
\end{proof}
\end{annexitem}

\begin{annexitem}{annex:p7.cr.np-conp}{thm:p07-cr-np-conp}{Short tautology certificates and complementation}
\begin{proof}
The tautology language is nonempty. The preceding theorem makes
polynomially bounded proof-system existence equivalent to
\(\mathsf{TAUT}\in\mathsf{NP}\). Non-tautology is in NP:
on a well-formed formula guess values for its finitely many
occurring atoms and evaluate its parse tree; accept a falsifying
assignment. Malformed codes, if excluded from TAUT, are checked
and accepted in polynomial time by this complement recognizer.
Hence TAUT is in coNP.

If NP=coNP, TAUT is therefore in NP. Conversely assume TAUT is
in NP and let \(L\in\mathsf{NP}\). By the Cook--Levin
reduction named in the body, construct in polynomial time a
well-formed formula \(H_x\) satisfiable exactly when \(x\in L\).
Then \(\neg H_x\) is a tautology exactly when \(x\notin L\).
Computing that formula and applying the assumed NP tautology
recognizer recognizes the complement of L in NP. Thus NP is
closed under complementation. If \(K\in\mathsf{coNP}\), its
complement is in NP and applying that closure puts K in NP,
so coNP is contained in NP. Conversely, for \(K\in\mathsf{NP}\)
its complement is in NP by closure, so K is in coNP. The two
classes are equal. This is the complete reduction argument after
the separately stated Cook--Levin premise, not a fresh proof of
that premise hidden inside a citation.

If P=NP, closure of deterministic polynomial time under complement
gives NP=coNP. The converse would additionally require deterministic
recognition of the NP certificate languages, which no step above
supplies. In particular, an NP tautology certificate system permits
nondeterministic guessing, not necessarily polynomial proof search.
\end{proof}
\end{annexitem}

\begin{annexitem}{annex:p7.cr.simulation-properties}{thm:p07-cr-simulation-properties}{Composition and the separate quantifier regimes}
\begin{proof}
Identity translates each system into itself. If
\(h(g(x))=f(x)\) and \(k(t(z))=h(z)\), then
\(k(t(g(x)))=f(x)\). The running time of g bounds its output
length by a polynomial; the running time of t on that output is
therefore polynomial in |x|. Thus composition is polynomial-time
and proves transitivity. Requiring translations in both directions
is consequently reflexive, symmetric and transitive.

If f is polynomially bounded by q and g has polynomial output
bound r, take any y in the common range and an f-proof x with
\(|x|\leq q(|y|)\). Then g(x) is an h-proof of y with length
at most \(r(q(|y|))\). This proves polynomial boundedness of h.
Choosing a minimum f-proof instead gives
\(s_h(y)\leq r(s_f(y))\). These existential choices are used
only to establish inequalities. A given shortest-proof inequality
instead has the quantifiers ``for every y there is some short
h-proof.'' It contains no machine description, time bound or
identity on every input proof. One may enumerate target proofs
until finding y, but that search is not licensed a polynomial
time bound by the inequality. This distinguishes the two notions
without claiming an unproved complexity-class separation between them.
\end{proof}
\end{annexitem}

\begin{annexitem}{annex:p7.cr.frege-comparison}{thm:p07-cr-frege-comparison}{Fixed finite rule replacements and all encoding cases}
\begin{proof}
First, simultaneous substitution preserves a Frege derivation.
At a hypothesis line substitute its formula. If a line is an
instance \(\sigma D\) of a rule \((C_i)/D\), a further
substitution \(\tau\) makes it an instance with composed
substitution \(\tau\sigma\) and premises \(\tau\sigma C_i\).
The same argument includes zero-premise axioms. Induction on
the sequence proves the substitution lemma, preserving its line count.

For each rule R of \(\mathcal F_1\), implicational completeness
of \(\mathcal F_2\) supplies a fixed derivation \(\pi_R\)
of its conclusion from its premises. There are finitely many R,
so fix those derivations once. For a source proof line produced
by R and substitution \(\sigma\), splice in \(\sigma\pi_R\),
referring its hypothesis uses to the already translated earlier
premises. Hypothesis occurrences are not added as unproved new
assumptions. An original hypothesis line is retained. The substitution
lemma proves that every new block has the required last line,
so induction gives the same final conclusion from the same hypotheses.

The maximum fixed block length gives \(\lambda'\leq c\lambda\).
Every block formula contains only a fixed number of occurrences
of its rule variables. Substitution formulas for variables occurring
in the rule are subformulas of the actual premise or conclusion
instances, each bounded by \(\rho\). Any extra variable in the
fixed target derivation can be mapped to a fixed one-atom formula.
Thus \(\rho'\leq c\rho\). Parse the supplied annotations,
copy the fixed templates, instantiate and reindex their references.
All work and output bits are polynomial in the complete encoded
input proof. Fresh auxiliary identifiers can be allocated outside
its finitely many original identifiers using short binary indices.

If the source proof map sends an invalid code to its default
tautology, send it to a fixed target proof of that particular
tautology. Such a proof exists by completeness and is fixed once,
so it costs constant time. This also handles two different choices
of default tautology. Valid codes translate their literal same
conclusion. The transformation therefore proves the all-string
p-simulation identity, not just a statement about well-formed codes.
Swap the two systems for the reverse translation.
\end{proof}
\end{annexitem}

\begin{annexitem}{annex:p7.cr.nd-comparison}{thm:p07-cr-nd-comparison}{Ordered contexts, simultaneous discharge and finite-basis bridges}
\begin{proof}
Negation and disjunction can be represented without duplicating
their variable operands over any adequate basis of arity at most
two. Here is the finite Boolean reason behind that source convention.
Adequacy supplies constant true and false formulas in a fixed
dummy atom. Some connective is nonmonotone, or every definable
function would be monotone; fixing its other argument to a
constant along a decreasing truth-table edge produces negation.
Some binary connective is nonaffine, or closure under composition
would give only affine functions. Its two-input table has exactly
one or three true entries: the remaining six two-true-entry functions
and the constant tables are affine. With optional input and output
negations, a one-true-entry table represents conjunction, using
each input once. De Morgan's transformation gives disjunction,
again using each input once. All dummy constant formulas are
fixed, so their repeated occurrences cost a constant per operand.
This proves the required linear-length representations N and O.

Substituting a derivation in a source ND system and prefixing a
fixed context \(\Delta\) to every line preserves that derivation.
For a rule instance already prefixed by \(\Delta'\), its transformed
prefix is the concatenation of \(\Delta\) and the substituted
\(\Delta'\), so it remains a legal common-prefix instance.
Induction on its lines proves the ND counterpart of the Frege
substitution lemma. Replace each fixed rule of \(\mathcal N_1\)
by a fixed derivation in \(\mathcal N_2\), whose existence
follows from implicational completeness. Instantiate it with the
same substitution and common prefix. Each block has a fixed number
of lines; its formula and context count is a constant times that
of the original line and its premises, because added prefixes
occur a fixed number of times. The copying and reindexing algorithms
are polynomial in the encoded sequence. This proves the two
linear source-measure bounds between ND systems.

For an ND system \(\mathcal N\), construct the finite Frege
system \(\mathrm{fr}(\mathcal N)\) as follows. For each ND
rule \(\mathcal A/L\), use a fixed formula rule on the
representations \(\mathcal A^*/L^*\), and a second one on
\((P\mathcal A)^*/(PL)^*\), where P is a fresh atom
standing for a common disjunctive prefix. Both are semantically
sound: if the common disjunct is true the conclusion is true;
otherwise the original line-rule entailment applies. An actual
rule instance with empty prefix uses the first rule. With nonempty
prefix \(\Delta=(D_1,\ldots,D_t)\), substitute
\(\bigvee(\neg D_1,\ldots,\neg D_t)\) for P in the
second rule. Left association makes its result exactly the
representation of the prefixed source line. Hence replacing
each line of any ND derivation by its starred representation
is a Frege derivation of the same represented last line with
the same line count and only constant-factor l growth. This
also proves implicational completeness of \(\mathrm{fr}(\mathcal N)\):
interpret any valid formula consequence as one between empty-context
lines, use the completeness of \(\mathcal N\), and translate it.

Conversely, construct \(\mathrm{nd}(\mathcal F)\) with each
Frege rule acting on empty-context lines, and the two sound
ND rules that discharge the last hypothesis into
\(\neg P\lor Q\), or introduce it back from that disjunction.
Common prefixing permits their use under any earlier context.
To establish implicational completeness between arbitrary lines,
discharge each hypothesis of every given line, starting at the
right, to obtain its right-associated representation in an
empty-context line. Frege completeness derives the corresponding
representation of the desired conclusion; reverse discharge
reintroduces its ordered hypotheses. This gives the desired
line with exactly its original order, not a reordered set.
Every Frege proof embeds simply as empty-context lines.
The representation steps are polynomial because each context has
length at most its encoded input length; no entire nested earlier
derivation is repeatedly duplicated.

The fixed-system comparison proved above converts each constructed
system to any other Frege or ND system on the same basis.
Empty-context final lines preserve the literal proved formula.
As in the preceding annex item, invalid codes use a fixed proof
of the source default tautology. This completes the uniform
theoremhood p-simulations. The interpretation as reusable sequences
is essential: unfolding every reused subderivation into a tree
would change the measure being compared.
\end{proof}
\end{annexitem}

\begin{annexitem}{annex:p7.cr.extension-elimination}{thm:p07-cr-extension-elimination}{Conservative semantics and the exact line-only elimination bound}
\begin{proof}
Start with a truth assignment to all ordinary atoms of the hypotheses
and conclusion satisfying every hypothesis. Traverse the derivation.
For an extension \(P_i\leftrightarrow C_i\), the freshness condition
means its new atom has no assigned earlier role and is absent from
the fixed original hypotheses and conclusion. Its defining formula
contains only earlier atoms. Assign \(P_i\) the current truth value
of \(C_i\), making the new declaration true without changing any
earlier value. Ordinary Frege inferences preserve truth by soundness.
At the last line the original conclusion is true, proving extension
soundness. No free choice remains for a defined atom once its
earlier values have been fixed.

To eliminate definitions, first regard all
\(E_i=(P_i\leftrightarrow C_i)\) as additional Frege hypotheses.
Their dependency order makes recursive expansion well defined.
Substitute each defined atom by its completely expanded defining
formula, in reverse dependency order as composition of substitutions.
The Frege substitution lemma preserves all lines. It leaves the
original hypotheses and conclusion unchanged, since they have no
defined atoms. Each substituted definition becomes
\(D_i\leftrightarrow D_i\). A fixed Frege proof of
\(Q\leftrightarrow Q\), instantiated with \(D_i\), proves it
in a constant c lines. Prepend those m proofs and refer substituted
hypothesis occurrences to them. The ordinary proof has at most
\(\lambda(\pi)+cm\) lines. Recursive expansion may duplicate
its defining subformulas at every level, so the copied formulas
and output string can be exponentially long. This argument
controls lines only and supplies no polynomial encoded translation.
\end{proof}
\end{annexitem}

\begin{annexitem}{annex:p7.cr.gate-compression}{thm:p07-cr-gate-compression}{All three skeleton cases, limited gates and fresh names}
\begin{proof}
For any defining set, induct on its subformula dependency order.
Each connective definition fixes the proxy's value uniquely from
already fixed children, including a constant node. This proves
the semantic defining-set claim. There is one constant-size equation
per distinct nonatomic subformula record, so the complete defining
set has \(O(m(A))\) atom occurrences. This includes every unary
node. The proof skeleton below need only introduce the gate equations
required at its finitely many rule-pattern nodes; it does not
enumerate the full internal tree of each substituted matrix.

For the skeleton translation, associate a proxy to every actual
subformula record needed by a fixed rule template, using identity
on atoms. Rename source auxiliary atoms into an ordinary private
namespace, disjoint from target hypothesis gates and gate names.
Process the original proof sequence online. A proxy associated
with a substitution matrix may be used as an atom without unfolding
that matrix's internal tree. Only connective nodes occurring in a
fixed rule or extension template require defining equations at that
step. There are a fixed number of them per step. A proxy already
defined or supplied as a hypothesis is reused; a new gate is
introduced only after its immediate dependency proxies are available.
Each source conclusion \(B_i\) is represented by \(P_{B_i}\).

For a hypothesis \(B_i\in\operatorname{def}_\kappa(A)\),
its constant connective pattern expresses one target gate equation.
The corresponding equation is in \(\operatorname{def}_{\kappa'}(A)\).
The finite definitions evaluating that fixed pattern under the
target basis imply \(P_{B_i}\); its truth-table inference is
a fixed sound finite-template consequence, hence has a fixed
target Frege derivation by implicational completeness. Instantiate
it with the actual child and gate proxies.

For an ordinary rule \((C_1,\ldots,C_r)/D\), its finitely many
pattern-gate equations together with the proxy premises
\(P_{C_1},\ldots,P_{C_r}\) entail \(P_D\): any assignment
consistent with the equations evaluates the original sound rule.
Fix a derivation of this finite consequence in the target Frege
system. Instantiate its pattern proxies \(P_E\) with
\(P_{\sigma E}\), for the source instance substitution \(\sigma\).
The proxy premises are earlier translated conclusions, while all
its equation instances are available or are freshly introduced
at these fixed pattern nodes. No expansion of \(\sigma E\)
is required when E is a rule variable. This proves this skeleton
step with a constant number of constant-width formula lines.

For an extension line \(P\leftrightarrow C\), first make
available just the proxy operands needed to evaluate the fixed
representation of this biconditional. Then introduce
\(P\leftrightarrow P_C\). The source P is still fresh:
no earlier source formula, original hypothesis or final conclusion
contains it, and the online construction has not preintroduced
future formula gates involving it. After that declaration introduce
the finitely many connective-pattern gates computing the source
extension sentence and derive its proxy from their equations.
The source conclusion proxy is not a fresh extension declaration;
it is the represented last line. Existing hypothesis/conclusion
proxies are kept as original atoms. This gives the extension case
without violating either freshness or the final-formula exclusion.

All blocks use fixed finite patterns depending only on the two
rule sets and connective bases. Hence their line and width costs
are bounded by constants, giving \(\lambda'\leq c\lambda\)
and \(\rho'\leq d\). Computing subformula identities and
substitution instances takes polynomial time in the encoded
input derivation, even when its formulas are much larger than
their l-count because of unary chains or long atom identifiers.

Finally, to close defining hypotheses in the same basis, rename
every new auxiliary extension variable away from all subformulas
of A. Simultaneously replace each nonatomic hypothesis proxy
\(P_E\) by E, leaving ordinary atoms alone. The conclusion
\(P_A\) becomes A, each hypothesis becomes
\(E\leftrightarrow E\), and target auxiliary declarations remain
fresh and acyclic. Prove each of these reflexive biconditionals
by a fixed theorem template and use the substitution lemma on
the rest. The line overhead is proportional to the number of
distinct subformula records of A; the maximum encoded formula
size is bounded by a constant times the earlier width times
the encoded size of A. These are the precise closing bounds
when unary constructors and atom encodings are counted.
\end{proof}
\end{annexitem}

\begin{annexitem}{annex:p7.cr.ef-line-compression}{thm:p07-cr-ef-line-compression}{The complete defining-set, skeleton and closing construction}
\begin{proof}
Rename auxiliary atoms so that the given proof's extension variables
and the proposed proxies for subformulas of A are distinct. Under
\(\operatorname{def}_{\kappa}(A)\), derive
\(P_E\leftrightarrow E\) in subformula order for each subformula E.
For an atom the equivalence is reflexive. For a constant use its
definition. For a unary or binary node, its definition and the
already obtained child equivalences give the desired equivalence
by a fixed sound congruence template. Implicational completeness
supplies a fixed finite derivation of each such template, depending
only on the connective, so this induction has \(O(m(A))\) lines
and maximum l-count \(O(l(A))\). At the root it gives
\(P_A\leftrightarrow A\).

Run the original proof under those additional hypotheses and
combine its last line A with this equivalence to derive \(P_A\).
Its original extensions remain legal: their variables occur in
neither A nor any of its defining-set hypotheses, and were renamed
away from all proxy identifiers. The resulting derivation has
\(\lambda\leq\lambda(\pi)+c m(A)\). Apply the preceding
skeleton lemma with the same source and target basis; the output
has a constant maximum l-count and only a constant-factor increase
in lines. Close its defining hypotheses by replacing subformula
proxies with the actual subformulas of A. The closing part of
that lemma adds \(O(m(A))\) lines and has maximum l-count
\(O(l(A))\). This proves the displayed two bounds for the
same original conclusion A.

Every stage is effective: parse the original proof and formula,
assign short fresh gate names, instantiate fixed finite derivations,
and make the indicated simultaneous substitutions. A subformula
substitution replaces an atom by a string of at most |A| bits.
The number of original records and every parsing operation are
polynomial in the supplied encoded input. Thus the construction's
running time and output string are polynomial, not merely a
nonuniform short-proof bound. Counting all constructors in m
keeps the line overhead valid even for long unary chains.
\end{proof}
\end{annexitem}

\begin{annexitem}{annex:p7.cr.ef-relative-bounds}{thm:p07-cr-ef-relative-bounds}{Different bases, monotone length bounds and encoded polynomial boundedness}
\begin{proof}
Let B be a tautology over the target basis. The source defining
set \(\operatorname{def}_{\kappa}(B)\) semantically entails
\(P_B\): its gate equations evaluate exactly B under every
ordinary-atom assignment. It has \(O(m(B))\) total occurrences
and constructor nodes, with fixed connector representations.
Represent this entailment by the nested disjunction of the negated
defining hypotheses followed by \(P_B\). The read-once
negation/disjunction representations proved in the ND comparison
make a source tautology T with \(m(T)\leq c m(B)\).
The bound hypothesis supplies its source extension proof with at
most \(L(m(T))\leq L(c m(B))\) lines.

Using each defining hypothesis, eliminate the corresponding
outer disjunction by a fixed sound Frege inference. Their number
is \(O(m(B))\), and their complete ordinary atom set occurs
already in T. Rename the proof's new extension atoms away from
those hypotheses. This gives a legal derivation of \(P_B\)
from the defining set with at most \(c L(c m(B))\) lines,
since \(L(n)\geq n\) absorbs the extra elimination steps.
Apply the cross-basis skeleton lemma to this derivation, keeping
the conclusion proxy and changing the gate equations to the
target basis. The result has constant maximum l-count and the
same line bound up to a fixed factor. Close its target defining
hypotheses to get B, adding \(O(m(B))\) lines and increasing
the maximum l-count to \(c l(B)\). This proves the claimed
relative bounds. It asserts existence of a source proof of T
by the bound hypothesis, not a polynomial algorithm that finds
that proof from B.

Replacing L by \(L^\uparrow(n)=\max_{j\leq n}L(j)\) is valid
for every original length bound and preserves polynomial growth.
If one extension system has a polynomial line bound in m,
the displayed transfer gives the same property for every fixed
adequate basis. The short-line theorem then makes all their
proofs polynomial in encoded formula length: there are polynomially
many lines with polynomially many leaf occurrences and at most
a polynomial number of auxiliary atoms. Rename those auxiliary
atoms to unused short binary identifiers, while preserving the
original atoms of B and their given encodings. Every original
identifier has length at most |B|; every auxiliary identifier
can have \(O(\log(m(B)+\lambda+1))\) bits. The short-line
construction copies only subformula strings of B and fixed
patterns, so unary syntax contributes at most a further |B|
factor. Rule annotations and parent references are polynomial
as well. This gives a polynomial bound in |B| on the full proof code.

Conversely, a polynomial encoded bound for one fixed extension
system gives a polynomial line bound in m after canonical atom
renaming. For a formula B, rename its t distinct atoms to short
identifiers of at most \(O(\log(t+1))\) bits, with t no greater
than m(B). The renamed formula has \(O(m(B)\log(m(B)+2))\)
bits, so its encoded short proof has polynomially many lines.
Rename back the original atoms and use fresh disjoint names
for its auxiliary atoms; the number of lines is unchanged.
All rule schemata and the freshness conditions are invariant
under this injective renaming. Hence this is a polynomial line
bound in m(B). The source leaf measure l by itself would not
count arbitrary unary syntax, so m is explicitly used here.

An ordinary Frege proof is an extension proof with no extensions.
The ND-to-Frege transformation multiplies its line count by a
fixed constant. If either had polynomial line bounds, the
extension system would too, and the equivalence just proved
would give a polynomially bounded propositional proof system.
The NP/coNP theorem then gives NP=coNP. Its contrapositive
proves the corrected corollary with NP unequal to coNP, not
merely with P unequal to NP.
\end{proof}
\end{annexitem}

\begin{annexitem}{annex:p7.cr.substitution-simulation}{thm:p07-cr-substitution-simulation}{A shared guarded proof and reverse-order pruning}
\begin{proof}
Let the extension declarations in a valid proof be
\(E_i=(P_i\leftrightarrow C_i)\), in their order of introduction,
and let \(D_k=((E_1\land E_2)\land\cdots)\land E_k\).
If there are no extensions the original proof already belongs
to the substitution system, so assume k positive.
First treat the declarations as additional ordinary Frege hypotheses.
Construct an ordinary no-hypothesis Frege proof of \(D_k\to A\)
as follows. Produce all projections \(D_k\to E_i\) with
\(O(k)\) lines: from the root conjunction derive the implication
to its two children, then share the prefix implications when
deriving the next two children, using transitivity of implication.
No separate k-step projection derivation is repeated for every leaf.

Guard every original ordinary proof line \(B_i\) by
\(D_k\to B_i\). For a declaration use its projection. For an
axiom use its ordinary axiom instance and the fixed tautological
derivation of \(B_i\to(D_k\to B_i)\). For a rule
\((C_1,\ldots,C_r)/E\), the fixed guarded consequence
\(D_k\to C_1,\ldots,D_k\to C_r\models D_k\to E\)
has a fixed Frege derivation: it is valid because assignments
falsifying the guard satisfy the conclusion, and assignments
satisfying it reduce to the original sound inference.
Instantiate that finite template at each step. This gives
\(D_k\to A\) in \(O(\lambda(\pi)+k)\) lines, with
maximum l-count \(O(k\rho(\pi)+\rho(\pi))\).

Now apply substitution only to the current no-premise theorem
\(D_k\to A\), replacing \(P_k\) by \(C_k\).
Freshness ensures \(P_k\) occurs in no earlier \(E_i\),
and is absent from A and from \(C_k\). Only the last guard
conjunct changes, becoming \(C_k\leftrightarrow C_k\).
Prove that reflexive biconditional by a fixed instance proof;
the current implication and it yield \(D_{k-1}\to A\)
by a fixed propositional template. Repeat with \(P_{k-1}\),
then down to \(P_1\). Each expanded self-biconditional is
removed immediately, so earlier definitions are not recursively
expanded. The final empty guard yields A itself. Each pruning
costs one substitution and a fixed number of Frege lines.
Since \(k\leq\lambda(\pi)\) and \(\rho(\pi)\geq1\),
the line count and maximum l-count are both bounded by
\(c\lambda(\pi)\rho(\pi)\), as stated.

The procedure parses the supplied proof, copies its definitions
and fixed guarded templates, and substitutes matrices already
present in the input. Guarded formulas and their annotations
have polynomial encoded length; the pruning never expands the
whole definition circuit. Thus this is a polynomial-time
translation of proof strings, not only an existence argument.
Invalid extension-proof codes are sent to a fixed substitution-Frege
proof of the source default tautology. This makes the identity
hold on every string. All substitution steps are theorem-only;
the process never substitutes into an undischarged premise.
\end{proof}
\end{annexitem}

\begin{annexitem}{annex:p7.cr.php-compression}{thm:p07-cr-php-compression}{Every smaller-instance clause, source counts and the final conclusion}
\begin{proof}
Let the clauses of \(\mathcal P_n\) hold under an assignment.
For each \(i<n\), if some \(P_{ij}\) with j<n-1 is true,
then the corresponding \(B_{ij}\) is true. Otherwise the
i-row forces \(P_{i,n-1}\). Its collision clause with the
nth pigeon makes \(P_{n,n-1}\) false. The nth row then forces
some \(P_{nj}\) with j<n-1, making that \(B_{ij}\) true.
This proves every smaller row disjunction, including its
complete last-hole case.

For \(i<k<n\) and j<n-1, suppose both \(B_{ij}\) and
\(B_{kj}\) hold. If both use their first disjuncts, then
\(P_{ij},P_{kj}\) violate the original j-hole collision.
If only the first uses its first disjunct, then
\(P_{ij},P_{nj}\) violate the collision between i and n.
The reverse mixed case violates the collision between k and n.
If both use their second disjuncts, then
\(P_{i,n-1},P_{k,n-1}\) violate the last-hole collision.
These four cases prove every smaller collision clause.
They do not assume functionality clauses for one pigeon's
different hole variables.

These finite Boolean arguments are uniform Frege derivations.
One can choose a convenient fixed sound complete basis containing
the fixed case-split, disjunction manipulation and collision
templates; each is itself a fixed consequence with a fixed
derivation in any other source Frege basis. A smaller row is
derived by eliminating the \(n\)-term \(i\)-row. Its early-hole cases
directly introduce a \(B_{ij}\); the last-hole case uses the
nth row and the single last-hole collision to eliminate its
\(n-1\) choices. Introducing and reassociating the \(O(n)\)-term
disjunctions takes \(O(n)\) fixed steps per choice, so a row needs
at most \(O(n^2)\) lines with \(O(n)\) leaf occurrences each.
There are \(n-1\) such rows. Each smaller collision uses the four
fixed two-disjunct cases just listed and at most four original
collision inputs, with a constant-size conclusion pattern.
There are \(O(n^3)\) collisions. Thus one induction step has
\(O(n^3)\) lines and \(O(n)\) maximum leaf count.

At \(n=2\) the clauses are \(P_{11},P_{21},
\neg P_{11}\lor\neg P_{21}\), which derive contradiction
by a fixed Frege inference. Compose the smaller-instance
substitutions and apply the substitution lemma to each subsequent
step. Summing \(O(m^3)\) over \(m=2\) through \(n\) gives \(O(n^4)\) lines.
Without abbreviations, every substitution replaces one atom
by the three-leaf \(B_{ij}\) tree; a \(k\)-stage fully expanded
proxy formula has \(3^k\) leaf occurrences before any
optional simplification. This proves exponential growth of
that unshared construction, not a minimum Frege lower bound.

For the shared construction, let \(Q^0_{ij}=P_{ij}\).
At stage \(k\) with \(m=n-k\) introduce fresh variables
\[
 Q^{k+1}_{ij}\leftrightarrow
 \bigl(Q^k_{ij}\lor(Q^k_{i,m-1}\land Q^k_{mj})\bigr),
 \qquad i<m,\ j<m-1.
\]
Definitions depend only on the preceding stage. The earlier
four collision cases and row argument apply with these proxy
variables and their definitions. There are \(O(n^3)\) definitions
over all stages; all are constant-size patterns. The derivation
from the original individual clauses still has \(O(n^4)\) lines
and \(O(n)\) maximum leaf count, with no accumulated substitution
tree in any line.

To turn this into a theorem proof with the claimed total size,
balance the original row disjunctions and the conjunction \(F\)
of all original clauses. Its tree has \(N=\Theta(n^3)\) leaf
occurrences and depth \(O(\log(n+2))\). Introduce the ordinary
acyclic gate definitions for \(F\), with final proxy \(H\). From \(H\)
and these gate equations derive each original clause, sharing
the conjunction projections. Guard the preceding pigeonhole
derivation by \(H\), replacing each ordinary inference by its fixed
guarded template. Its own fresh proxy definitions may be added
unconditionally and then guarded by a fixed tautological step.
The resulting no-hypothesis proof has \(O(n^4)\) internal lines
of \(O(n)\) leaves and ends in \(\neg H\).

Derive \(H\leftrightarrow F\) from the balanced gate equations
by ordinary congruence in subformula order. A node's derivation
uses a fixed number of formulas whose size is a constant times
its expanded subtree. Summed subtree leaf counts are
\(O(N\log(n+2))\), because each leaf occurs in at most that many
ancestor subtrees. This bridge therefore contributes only
\(O(N\log(n+2))\) total leaf occurrences; its final largest formula
has \(O(N)\) leaves. Combining it with \(\neg H\) gives
\(\neg F=A_n\), with no new atom in that final conclusion.
All original atom, stage proxy and gate names have \(O(\log(n+2))\)
bits, as do line references in an \(O(n^4)\)-line sequence. Total
encoded proof length is consequently \(O(n^5\log(n+2))\); the
bridge is smaller than that bound. The largest full-proof line
must include the \(\Theta(n^3)\)-leaf final formula, so the narrower
\(O(n)\) width bound was correctly confined to the clause derivation.
\end{proof}
\end{annexitem}

\begin{annexitem}{annex:p7.cr.small-php}{thm:p07-cr-small-php}{All nineteen clauses and the two duplicated inputs}
\begin{proof}
Write \(x_i=P_{i1}\) and \(y_i=P_{i2}\).
The following table gives the complete refutation on
Sheet~\ref{sheet:resolution-cpl}. Clauses are sets, so
duplicate literals in a resolvent are identified. All input
nodes are included. A numbered pair records the earlier parents
and the pivot is the atom resolved between them.
\begin{center}\small
\begin{tabular}{@{}rlll@{}}
\toprule
Node & Clause & Parents & Pivot \\
\midrule
1 & \(x_1\lor y_1\) & input & -- \\
2 & \(x_2\lor y_2\) & input & -- \\
3 & \(x_3\lor y_3\) & input & -- \\
4 & \(\neg x_1\lor\neg x_3\) & input & -- \\
5 & \(\neg x_2\lor\neg x_3\) & input & -- \\
6 & \(\neg y_1\lor\neg y_2\) & input & -- \\
7 & \(\neg y_1\lor\neg y_3\) & input & -- \\
8 & \(\neg y_2\lor\neg y_3\) & input & -- \\
9 & \(\neg x_1\lor\neg x_2\) & input & -- \\
10 & \(y_1\lor\neg x_3\) & 1,4 & \(x_1\) \\
11 & \(y_2\lor\neg x_3\) & 2,5 & \(x_2\) \\
12 & \(\neg x_3\lor\neg y_2\) & 10,6 & \(y_1\) \\
13 & \(\neg x_3\) & 12,11 & \(y_2\) \\
14 & \(x_1\lor\neg y_3\) & 1,7 & \(y_1\) \\
15 & \(x_2\lor\neg y_3\) & 2,8 & \(y_2\) \\
16 & \(\neg y_3\lor\neg x_2\) & 14,9 & \(x_1\) \\
17 & \(\neg y_3\) & 16,15 & \(x_2\) \\
18 & \(y_3\) & 3,13 & \(x_3\) \\
19 & empty clause & 18,17 & \(y_3\) \\
\bottomrule
\end{tabular}
\end{center}
Each row removes exactly the listed complementary pivot from
its two parents and unions their remaining literals, so the
table checks every inference directly. Nodes 1 and 2 each have
two later uses, once in the proof of \(\neg x_3\) and once
in the proof of \(\neg y_3\). Every other input is used once
and every derived clause has one subsequent use. The DAG has
nine input nodes and ten derived nodes. Its binary-tree
unfolding duplicates exactly nodes 1 and 2, producing eleven
input occurrences and the same ten internal occurrences,
hence 21 total clause occurrences. Comparing those two explicit
representations does not claim either is a minimum refutation.
Its conclusion is unsatisfiability of \(\mathcal P_3\), or
equivalently tautology of \(A_3\), not a general resolution
lower bound for the larger family.
\end{proof}
\end{annexitem}
